\chapter{Drooling (Sialorrhea)}

\section*{Definition}
The unintentional loss of saliva from the mouth, resulting from excessive saliva production (hypersalivation) or impaired oral motor control causing difficulty swallowing saliva, which can lead to skin irritation, aspiration, and social embarrassment.

\section*{What Happens in Your Body}
Saliva is produced by three major paired salivary glands (parotid, submandibular, sublingual) and minor glands, with daily production of 0.5-1.5 L. Drooling occurs due to excess production (rare) or, more commonly, from oropharyngeal dysphagia, impaired lip seal, poor head control, or reduced swallowing frequency. Neurological conditions disrupt the coordinated neuromuscular sequence of saliva clearance.

\section*{Causes}
\begin{itemize}
  \item Neurological disorders (cerebral palsy, Parkinson disease, ALS, stroke)
  \item Developmental delay (children with poor oral motor control)
  \item Oropharyngeal dysphagia (difficulty swallowing)
  \item Dental problems (teething in infants, ill-fitting dentures)
  \item GERD (increased salivation as compensatory response)
  \item Medications (clozapine, risperidone, cholinesterase inhibitors, ketamine)
  \item Toxins (mercury, organophosphates, arsenic)
  \item Pregnancy (hormonal changes causing hypersalivation)
  \item Mouth breathing (nasal obstruction, enlarged tonsils/adenoids)
  \item Parkinson disease (stooped posture, reduced swallowing frequency, rigidity)
\end{itemize}

\section*{When to See a Doctor}
See your doctor if:
\begin{itemize}
  \item Symptoms are severe or getting worse over time
  \item Sudden onset of drooling (possible stroke, especially with facial droop or speech difficulty), drooling with fever and difficulty swallowing (possible epiglottitis or retropharyngeal abscess)
  \item You have other concerning symptoms such as fever, weight loss, or bleeding
\end{itemize}

\section*{Related Symptoms}
\begin{itemize}
  \item Dysphagia
  \item Speech difficulty
  \item Facial droop
  \item Bad breath
  \item Choking
\end{itemize}
\textbf{Important:} If this symptom persists, your doctor will check for serious underlying causes.

\section*{Self-Care / Home Management}
\begin{itemize}
  \item Practice good oral hygiene and use mouth rinses to reduce saliva pooling and odor.
  \item Swallow frequently and consciously, especially when concentrating or in a relaxed posture.
  \item Use absorbent tissues or a small towel to manage visible drooling in social settings.
  \item Stay hydrated to avoid thick, tenacious saliva that is harder to swallow.
  \item Consult a dentist for dental causes such as ill-fitting dentures or oral infections.
\end{itemize}

\section*{How Your Doctor Will Evaluate You}
\begin{itemize}
  \item Your doctor will assess oral motor function, swallowing, and neurological status.
  \item A neurological examination evaluates cranial nerves, muscle tone, and coordination.
  \item A swallow study (videofluoroscopic swallow study) may be ordered to assess aspiration risk.
  \item Medication review helps identify drugs that may cause or worsen drooling.
  \item Referral to a neurologist, speech-language pathologist, or ENT specialist may be arranged.
\end{itemize}

\section*{Treatment Options}
\begin{itemize}
  \item Treatment is directed at the underlying neurological or structural cause.
  \item Anticholinergic medications (glycopyrrolate, scopolamine patch) reduce saliva production.
  \item Botulinum toxin injections into salivary glands for refractory cases.
  \item Speech therapy and oral motor exercises improve swallowing efficiency.
  \item Surgical options (duct ligation, salivary gland excision) in severe, treatment-resistant cases.
\end{itemize}

\section*{References}
\begin{thebibliography}{99}
\bibitem{ref1} Hockstein NG, Samadi DS, Gendron K, et al. "Sialorrhea: a management challenge." Am Fam Physician. 2004;69(11):2628-2634.
\bibitem{ref2} McGeachan AJ, McDermott CJ. "Management of sialorrhoea in motor neurone disease." J Neurol Neurosurg Psychiatry. 2017;88(1):85-90.
\bibitem{ref3} Lakraj AA, Moghimi N, Jabbari B. "Sialorrhea: anatomy, pathophysiology and treatment." Clin Neuropharmacol. 2013;36(1):28-36.
\bibitem{ref4} Squires N, Wills A, Rowson J. "The management of drooling in adults with neurological conditions." J Neurol. 2020;267(12):3713-3720.
\bibitem{ref5} Blasco PA, Allaire JH. "Drooling in the developmentally disabled: management practices and recommendations." Dev Med Child Neurol. 1992;34(10):849-862.
\bibitem{ref6} Roberts A. "Sialorrhea in adults: treatment." UpToDate. Wolters Kluwer; 2025.
\end{thebibliography}

\clearpage